\chapter{Capstone 项目工作流、项目模板与可信交付}

\section{案例导读：Capstone 不是最后一份作业，而是一次可信交付}

前五部分训练的是单项能力：会配置环境、会处理数据、会建立 baseline、会评估模型、会把现代 AI 工具放进验证边界。Capstone 要做的事情不同：它要求学生把这些能力组织成一份别人能够复跑、复查、评分和继续维护的项目包。

本章是 Part VI 的入口。请不要把它读成普通项目管理清单，而要把 project board 看成一种科研证据结构：每一项任务都要说明它服务哪个科学问题、留下什么证据、遇到边界时如何停下来，以及最终由谁签字负责。

\section{案例目标}

走到这里，学生通常已经学过 Python、数据处理、机器学习、案例分析，以及 Part V 里一整条“现代 AI 进入科研 workflow 之后怎样保持可信”的主线。真正临门一脚的问题就变成了：\textbf{怎样把这些分散能力收束成一个能展示、能复跑、能签字负责的期末项目。}

本章的目标有四点：

\begin{enumerate}
    \item 理解 Capstone 项目不只是“再训练一个模型”，而是一份要同时满足科学问题、数据边界、验证证据、文献支撑和交付规范的综合工作流；
    \item 学会把项目状态写成一张结构化 project board，而不是把所有进度都留在口头记忆里；
    \item 学会把第 35 章到第 38 章的四类边界，收束到同一张项目卡上；
    \item 学会把 project board 转成 ready、review、revise 和 blocked 四类出口，以及更适合执行的 next-step checklist。
\end{enumerate}

\section{为什么 Part VI 不是“再加一个模型”}

很多学生一听到 capstone，就会本能地把注意力放到“我要训练什么更酷的模型”。但一个真实项目最后能不能站得住，往往不主要取决于模型名字，而取决于下面这些更朴素的问题：

\begin{itemize}
    \item 科学问题到底是否清楚，还是只有一个模糊题目；
    \item 数据能不能合法使用、能不能稳定复跑；
    \item 有没有一个值得比较的 baseline；
    \item 验证、误差分析和人工复核出口是否已经准备好；
    \item related work、claim ledger 和 caveat 是否已经整理；
    \item AI 到底参与了哪些步骤，哪些判断仍然必须由人签字。
\end{itemize}

如果这些问题都没有被写清楚，那么项目很容易出现一种表面上很忙、实际上很脆弱的状态：

\begin{itemize}
    \item notebook 里有很多单元格，但不知道哪一格是最终证据；
    \item 模型分数看起来不错，但没有 baseline、没有 diagnostics；
    \item 报告已经开始写了，但引用和 caveat 还没有整理；
    \item AI 确实帮了很多忙，但最终没有留下可检查的 AI Usage Log。
\end{itemize}

所以，本章的核心立场是：\textbf{Capstone 不是一个“更大模型”的章节，而是一个把整本书前面内容接成可信交付的章节。}

\section{一个极小的 capstone project board 数据集}

配套 notebook 使用 \nolinkurl{capstone_project_workflow_demo.csv}。这不是成绩册，也不是完整项目管理软件导出的 board，而是一组完全本地、教学用途的\textbf{capstone 项目卡片}。每一行代表一个学生项目的当前状态，包含：

\begin{itemize}
    \item \texttt{project\_id}；
    \item \texttt{theme}；
    \item \texttt{science\_question\_status}；
    \item \texttt{data\_license\_status}；
    \item \texttt{baseline\_status}；
    \item \texttt{validation\_status}；
    \item \texttt{literature\_ledger\_status}；
    \item \texttt{ai\_use\_statement\_status}；
    \item \texttt{report\_outline\_status}；
    \item \texttt{sharing\_boundary}；
    \item \texttt{model\_hype\_score}；
    \item \texttt{reference\_route}。
\end{itemize}

这里的 \texttt{reference\_route} 不是机器学习标签意义上的“真理类别”，而是教学上预设的人类参考路由，分成四类：

\begin{itemize}
    \item \texttt{ready\_for\_capstone}：项目已经具备进入正式执行与展示准备的最小条件；
    \item \texttt{manual\_review}：项目方向没有问题，但还有验证、文献或声明边界需要人工收口；
    \item \texttt{revise\_plan}：科学问题或 baseline 还没有站稳，需要回到计划层改写；
    \item \texttt{blocked}：当前存在数据许可或共享边界问题，不该继续往下推进。
\end{itemize}

这个数据集故意放入了几类最容易让学生在 capstone 阶段失速的项目：

\begin{itemize}
    \item \texttt{model\_hype\_score} 很高，但 baseline 仍然缺失的项目；
    \item 使用受限数据、目前还不能安全共享或复现的项目；
    \item 代码和实验主体大致齐了，但 literature ledger 与 AI-use statement 还没收口的项目；
    \item 其实已经很扎实，却因为“不够炫”而容易被忽视的项目。
\end{itemize}

因此，本章真正要训练的不是“怎样挑一个最酷的题目”，而是更重要的能力：

\begin{itemize}
    \item 你能不能判断一个项目是不是已经准备好进入 Capstone 节奏；
    \item 你能不能明确说出它为什么应该 ready、为什么应该 review、或者为什么现在必须停下来。
\end{itemize}

\section{先看一个很常见的 baseline：按模型兴奋度排优先级}

如果没有结构化 project board，一个很常见的 baseline 就是：\textbf{按 \texttt{model\_hype\_score} 从高到低排序，然后优先推进最“像 AI 项目”的题目。}

\begin{examplebox}[按 hype 分数排优先级的 baseline]
\begin{lstlisting}[style=pythonstyle]
top_projects = sorted(
    rows,
    key=lambda row: row["model_hype_score"],
    reverse=True,
)[:3]
\end{lstlisting}
\end{examplebox}

这种做法的问题非常直接：它把“看起来先进”误当成“已经准备好交付”，却没有先问：

\begin{itemize}
    \item 数据许可和共享边界是否清楚；
    \item baseline 和验证链条是否已经存在；
    \item literature ledger 是否已经补齐；
    \item 报告和 AI-use statement 是否已经能支持最终展示。
\end{itemize}

在本章教学数据上，按 \texttt{model\_hype\_score} 排名前三的项目里，真正已经可以直接进入 capstone 执行阶段的只有一个。因此 notebook 中会得到表~\ref{tab:ch39_baseline_vs_capstone} 这个最小对照：

\begin{table}[htbp]
\centering
\small
\setlength{\tabcolsep}{4pt}
\begin{tabular}{@{}p{0.23\textwidth}>{\centering\arraybackslash}p{0.16\textwidth}>{\centering\arraybackslash}p{0.18\textwidth}>{\centering\arraybackslash}p{0.15\textwidth}p{0.20\textwidth}@{}}
\toprule
方法 & 直接 ready 精度 & blocked 混入率 & review 出口 & 教学解读 \\
\midrule
只按 \texttt{model\_hype\_score} 取前三项 & 0.33 & 0.00 & 否 & 容易把计划未稳的高 hype 项目误当成优先交付对象 \\
结构化 capstone workflow & 1.00 & 0.00 & 是 & 先做项目级 gate，再决定哪些项目已经 ready \\
\bottomrule
\end{tabular}
\caption{本章 notebook 中两个最小工作流的对照。重点不是“谁更前沿”，而是谁更接近可签字交付。}
\label{tab:ch39_baseline_vs_capstone}
\end{table}

\section{capstone workflow 的关键，是把前四章的边界接到同一张项目卡上}

本章和前面四章的关系非常直接。一个项目卡之所以能支持 ready / review / blocked 的判断，是因为它同时继承了四类约束：

\begin{itemize}
    \item \textbf{第 35 章：代码与 notebook 验证边界}。也就是 \texttt{validation\_status} 不能只表示“代码跑过”，而要表示关键数值检查、回归检查和复跑证据是否已经准备好。
    \item \textbf{第 36 章：agentic workflow 的人工接管边界}。也就是项目不必什么都自动拍板，边界样本和边界决策应该允许进入 \texttt{manual\_review}。
    \item \textbf{第 37 章：文献阅读与 claim ledger 边界}。也就是 \texttt{literature\_ledger\_status} 要告诉你 related work、结果 claim 和 caveat 是否已经可回查。
    \item \textbf{第 38 章：AI 使用、版权、许可与声明边界}。也就是共享边界和 AI 使用说明都要写清楚：哪些材料可以共享、哪些步骤用了 AI、哪些结论仍由人签字。
\end{itemize}

如果项目使用了第 30--34 章的深度模型，或者第 35--38 章的 LLM / agent 工具，那么 Part V 收束章里的“可信交付包”也应该成为 project board 的一部分。深度模型至少要留下 Model Experiment Record、evaluation artifact、error / review artifact 和 Trust Statement；LLM 或 agent workflow 至少要留下 Evidence Record、AI Usage Log、verification artifact 和 disclosure note。

因此，本章 notebook 的最小 routing 逻辑不是“哪个题目更火”，而更接近下面这种结构：

\begin{examplebox}[一个最小 capstone routing 逻辑]
\begin{lstlisting}[style=pythonstyle]
if data_license_status != "clear":
    blocked
elif sharing_boundary != "local_only":
    blocked
elif science_question_status != "clear":
    revise_plan
elif baseline_status != "ready":
    revise_plan
elif validation_status != "ready":
    manual_review
elif literature_ledger_status != "ready":
    manual_review
elif ai_use_statement_status != "drafted":
    manual_review
elif report_outline_status != "ready":
    manual_review
else:
    ready_for_capstone
\end{lstlisting}
\end{examplebox}

请注意，这里的重点不是把项目管理“自动化得更厉害”，而是把那些最容易在期末阶段失控的边界显式写出来。

\section{为什么 \texttt{manual\_review} 不是失败出口}

很多学生在项目末期会感到焦虑，仿佛一个项目只要还没有达到 \texttt{ready\_for\_capstone}，就说明自己做得不好。其实在真实教学和科研里，这种理解经常是错的。

如果一个项目满足下面任一情况：

\begin{itemize}
    \item 核心实验已经大致成立，但 validation 记录还没写清；
    \item 结果段已经能成形，但引用和 caveat 还没有完全补齐；
    \item AI 的参与方式已经很多，但 usage statement 还没和实际记录对齐；
    \item 报告结构已经成形，但最后的 human signoff 还没有完成；
\end{itemize}

那么把它送进 \texttt{manual\_review} 往往不是“项目失败”，而是\textbf{项目管理正确}。

同样，\texttt{revise\_plan} 也不是“退步”。如果科学问题还很宽、baseline 仍然缺失，那么继续堆模型通常只会把后面的解释和展示做得更乱。早一点在 project board 上暴露这个问题，反而比晚一点在答辩前临时发现要安全得多。

\section{配套 notebook 做了什么}

本章 notebook 采用的是一个 teaching-first 的最小设计：

\begin{itemize}
    \item 先读取 project board，并打印 science question、data license、sharing boundary 和 reference route 统计；
    \item 用按 \texttt{model\_hype\_score} 排名前三的方式做一个 naive priority baseline；
    \item 再实现一个最小 capstone routing workflow；
    \item 输出 ready queue、manual review queue、revise-plan queue 和 blocked queue；
    \item 为每个非 ready 项目生成 next-step checklist；
    \item 最后打印一个更适合 capstone assistant 的 prompt 模板，以及 project snapshot。
\end{itemize}

这份 notebook 不依赖真正的项目管理服务，也不连接在线代理框架。它真正要教的是：\textbf{项目管理里最值得保留的，不是某个工具按钮，而是项目边界与交付证据的结构。}

\section{把路线图变成一个可执行项目模板}

如果把本章和整本书前面的主线压缩成一个最小 capstone 模板，那么一个真正可交付的项目至少应该能回答下面七件事：

\begin{enumerate}
    \item \textbf{科学问题}：你究竟想回答什么问题，为什么它值得回答；
    \item \textbf{数据与许可证}：数据来自哪里，能否公开重跑，哪些材料不能外发；
    \item \textbf{baseline}：在更复杂模型之前，有没有一个简单但可信的参考方案；
    \item \textbf{模型、AI 工具与诊断}：至少两个可比较实现，以及必要的误差分析、不确定度、异常样本解释或 Part V 可信交付包；
    \item \textbf{literature ledger}：related work、核心 claim 和 caveat 是否已经整理成可回查结构；
    \item \textbf{AI-use statement}：AI 具体帮助了哪一步，哪些判断没有交给 AI，哪些材料没有被发送到外部服务；
    \item \textbf{最终交付}：report、notebook、图表、Git 记录和展示结构是否已经准备好。
\end{enumerate}

也就是说，Capstone 模板不应该只是一页题目列表，而更像一份\textbf{项目签字前的最小证据清单}。

\begin{warnbox}[Capstone 阶段的常见误区]
\begin{itemize}
    \item 把“模型更复杂”误当成“项目更成熟”；
    \item 直到最后一周才想起 baseline、diagnostics 和 report outline；
    \item 代码、文献、声明和展示分别保存在不同地方，彼此没有连接；
    \item 只记得“用过 AI”，却说不清 AI 到底参与了哪些步骤；
    \item 发现共享边界或许可证问题后，仍然继续往外部服务粘贴材料。
\end{itemize}
\end{warnbox}

\section{和整本书的关系}

本章可以看成整本书的一次收束：

\begin{itemize}
    \item Part 0 到 Part II 提供 Python、数据和科学计算基础；
    \item Part III 提供机器学习 workflow、baseline、评估和可信度判断；
    \item Part IV 把这些工具放回具体天文与物理案例；
    \item Part V 进一步讨论深度学习与现代 AI 怎样进入科研 workflow；
    \item 本章则把这些能力收束成一个可展示、可复现、可签字的项目模板。
\end{itemize}

因此，真正的 Capstone 不是某一章之后突然冒出来的新任务，而是整本书前面所有能力的一次联合检查。

\section{AI 助手如何帮助你}

在这个案例里，AI 助手最适合帮助你：

\begin{itemize}
    \item 把分散的 notebook 进度、文献摘录和使用记录整理成 project board；
    \item 帮你把每个非 ready 项目的缺口改写成更清楚的 next-step checklist；
    \item 帮你把 literature ledger、AI-use statement 和 report outline 接成一个更清楚的项目骨架；
    \item 帮你把口头计划整理成项目展示大纲；
    \item 帮你检查说明文字是否遗漏了数据来源、验证边界和 human signoff。
\end{itemize}

但最终哪些项目已经可以交付、哪些材料不能共享、哪些 claim 足以写进报告，以及哪些结论值得你自己签字承担，仍然必须由你本人判断。

\section{练习}

\begin{exercisebox}[基础练习]
把教学数据里一个已经 \texttt{ready\_for\_capstone} 的项目改成 \texttt{sharing\_boundary = external\_unknown}。重新运行 notebook，观察它为什么会被直接 route 到 \texttt{blocked}。
\end{exercisebox}

\begin{exercisebox}[进阶练习]
请再添加一条新的项目卡：科学问题清楚、数据公开、baseline 已完成，但 \texttt{literature\_ledger\_status} 和 \texttt{report\_outline\_status} 都还是 \texttt{partial}。预测它会落到哪个 route，并说明原因。
\end{exercisebox}

\begin{exercisebox}[开放练习]
结合你自己的课程题目，写一张最小 capstone 项目卡。至少包含：科学问题状态、数据许可状态、baseline 状态、validation 状态、literature ledger 状态，以及一条你愿意公开提交的 AI-use statement 说明。
\end{exercisebox}

\section{本章小结}

Capstone 最容易被误解成“期末再做一个大模型”。但更可靠的理解是：\textbf{它是一份项目级 workflow，把问题定义、数据边界、baseline、验证、文献证据、AI 使用声明和最终展示收束到同一套证据结构里。}

只要你能持续把这些边界写进 project board，并愿意让项目在 ready、review、revise 和 blocked 之间诚实流动，那么最后交出来的成果就更可能不仅“会跑”，而且真正可解释、可复现、可签字。
